\subsection{Dimension of a finitely generated algebra}

PROPOSITION 3. --- Let \(\rho: A \to B\) be a ring homomorphism. Set \(n = \sup_{\mathfrak{p} \in \operatorname{Spec}(A)} \dim(B \otimes_A \kappa(\mathfrak{p}))\) . We have the inequality

\[
\dim (\mathbf {B}) + 1 \leqslant (\dim (\mathbf {A}) + 1). (n + 1).
\]

We may assume \(\dim(A) \neq -\infty\) and \(n < +\infty\) . Let \(\mathfrak{q}_0 \subset \ldots \subset \mathfrak{q}_m\) be a chain of prime ideals of B; set \(\mathfrak{p}_i = \rho^{-1}(\mathfrak{q}_i)\) . The sequence of \(\mathfrak{p}_i\) is increasing, so the set of its values has cardinality \(\leqslant \dim(A) + 1\) . For each \(\mathfrak{p} \in \operatorname{Spec}(A)\) , the set of \(\mathfrak{q}_j\) such that \(\mathfrak{p}_j = \mathfrak{p}\) is a chain in the subset \(({}^{a}\rho)^{-1}(\mathfrak{p})\) of \(\operatorname{Spec}(B)\), and therefore has cardinality less than or equal to \(\dim(B \otimes_A \kappa(\mathfrak{p})) + 1\) (no. 1, lemma 1), and consequently to \(n^o 1\)\(n + 1\) . It follows that \(m + 1 \leqslant (\dim(A) + 1)(n + 1)\) , whence the proposition.

Remark 1. --- If the rings A and B are Noetherian, we shall see below (§ 3, no. 4, cor. 2 to prop. 7) that one has the inequality dim(B) ≤ dim(A) + n, stronger than that of prop. 3.

COROLLARY 1. --- Suppose that one has  and that there exists an integer n such that \textless for all {}. Then we have \_\textless{}.

COROLLARY 2. --- Let A be a ring and {}{} the ring of polynomials in one indeterminate with coefficients in A. We have:

\[
1 + \dim (\mathrm{A}) \leqslant \dim (\mathrm{B}) \leqslant 1 + 2 \dim (\mathrm{A}).
\]

The first inequality has already been proved (§ 1, no. 3, example 4). Let us prove the second. For every prime ideal  of A, the ring {, which is isomorphic to }{}, is principal and is not a field, hence is of dimension 1 (§ 1, no. 3, example 2), and the inequality follows from Proposition 3.

Remark 2. --- We shall see later (§ 3, no. 4, cor. 3 to prop. 7) that, if A is Noetherian, one has {}{}.

However, whatever the integers \(n\) and \(q\) with \(n + 1 \leqslant q \leqslant 2n + 1\) , there exists a ring A of dimension \(n\) such that \(\dim(A[X]) = q\) (see p. 84, exerc. 7).

COROLLARY 3. --- If A is of finite dimension, every non-zero A-algebra of finite type is of finite dimension.

Indeed, from Corollary 2, by induction on \(n\), it follows that the ring \(\mathbf{A}[\mathbf{T}_{1}, \ldots, \mathbf{T}_{n}]\) is of finite dimension if \(\mathbf{A}\) is of finite dimension; a fortiori, every nonzero quotient of \(\mathbf{A}[\mathbf{T}_{1}, \ldots, \mathbf{T}_{n}]\)\(\S\) is of finite dimension (§ 1, no. 3, Proposition 6).

